% \bodyRerun* are defined in generated/body_numbers.tex (paper preamble); fallbacks only for the --check stub.
\providecommand{\bodyRerunN}{\ai{bodyRerunN}}\providecommand{\bodyRerunIdentical}{\ai{bodyRerunIdentical}}\providecommand{\bodyRerunFlips}{\ai{bodyRerunFlips}}\providecommand{\bodyRerunFlipPct}{\ai{bodyRerunFlipPct}}\providecommand{\bodyRerunChangedPct}{\ai{bodyRerunChangedPct}}\providecommand{\bodyRerunMedDiff}{\ai{bodyRerunMedDiff}}
This appendix documents every configuration of the optimization study of \S\ref{sec:opt}: what was run, how each tested configuration was chosen, and what it cost. The optimizer is \jevopt{}, a Jev-specific layer over GEPA running inside the OMP agent harness. Three pre-registered rounds used the same <<block:opt.nTasks>> task configurations and the same test splits. The first round compares the \emph{paper's question} (the zero-shot question built from the paper's own label definitions, as a program) with \emph{GEPA text search} and \emph{GEPA program search}. Its protocol was fixed before its runs, and all its test scores were recorded before the protocols of the other two rounds. The second round is \emph{GEPA with demonstrations and calibration} (class calibration, demonstrations, length-limited GEPA and ensembles as components of one configuration family). Its components and its protocol were recorded after the first round's test scores were known. Its chosen configurations were fixed before any of its test calls. The third round is \emph{question optimization}, called feature-question search in the code. It asks many question variants in packed requests, and their answer probabilities serve as features for a small logistic regression trained on the fit slice. The variants include questions the optimizer model rewrote from the task definition in the form Jev answers best, of which the best is also tested as written (the \emph{rephrased question}). Its code and protocol, with the selection rule, were also recorded after the first round's test scores were known. Its chosen configurations were fixed and tested before the second round's test. The second and third rounds were developed in parallel on validation. Each configuration was fixed before its own test runs and tested once. Table~\ref{tab:opt-prereg} gives the times. Alongside them ran a validation-only study of \emph{Jev-adapted GEPA}, a pre-registered rerun of the chosen configuration with two further optimizer seeds (\emph{seed variance}, Appendix~\ref{app:opt-seeds}), \emph{candidate reranking and fitted thresholds} for three large-label evaluated tasks, pre-registered \emph{calibration and logistic regression on the zero-shot evaluation's answers} for <<block:opt.headsN>> evaluated test sets (Appendix~\ref{app:opt-heads}), and a LifeTox pilot. The <<block:opt.nTested>> test evaluations of the optimizer-seed-0 runs are in Table~\ref{tab:opt-test-all}, those of the seed reruns in Table~\ref{tab:opt-seeds-headline}. Every number in this appendix is computed by \texttt{scripts/\allowbreak paper\_appendix\_opt.py} from the run directories under \texttt{runs/\allowbreak jev-opt/}, the analysis outputs under \texttt{results/} and \texttt{scripts/\allowbreak appendix/\allowbreak opt\_history.json}, which records, from the study's git history, when each protocol was recorded and the code constants of the \jevopt{} source in the version recorded by the runs that used them.

Two kinds of system appear side by side and are never merged. \emph{Zero-shot} rows run a program exactly as written. ``Jev + fitted numbers'' rows add class biases, Score cutoffs or thresholds fit on training data, and ``Jev + classifier on the answers'' rows use Jev's answer probabilities as features for a small logistic or ridge regression (a \emph{head}) trained on the labelled examples. Both are trained systems, closer to the papers' task-specific methods than to their zero-shot baselines. Validation scores of any configuration chosen on validation are optimistically biased. We state this bias and do not correct it.

\subsection{Tasks, data splits and the one-test-use ledger}\label{app:opt-data}
The tasks are the evaluated tasks with the largest gaps that a prompt might close: ConTempo TB-Dense, EZ-STANCE Subtask~A, SMARTER implicit hate, Re3-Sci edit intent, INLI, and DREsS\_New with its three rubrics as separate tasks (Table~\ref{tab:opt-data}). Train and validation data come only from each paper's released train and dev data (\texttt{jev-opt/\allowbreak experiments/\allowbreak README.md}): deterministic label-stratified samples, with the validation split set to the whole dev split where that is smaller (TB-Dense: the standard dev documents, DREsS: the reconstructed validation split). The test split is exactly the zero-shot evaluation's frozen English test set: each experiment's \texttt{prepare.py} maps every test item through the evaluation adapter's own state builder and asserts that the row count equals the zero-shot evaluation's $n$. Where the zero-shot evaluation screened test items for non-English text flag by flag, train and validation use an automatic approximation. INLI samples whole premises and expands each to all of its premise--hypothesis pairs, so no premise spans splits.
<<block:opt_data>>
From the second pre-registration on, all numbers and logistic regressions are fit on the first <<block:opt.fitExamples>> examples of a stratified ordering of train (the \emph{fit slice}). GEPA's reflection samples its minibatches from a separate pool of <<block:opt.reflPool>> train examples, so reflection never sees fit-slice examples except in aggregate. The hyperparameters of the logistic regression are chosen by <<block:opt.cvFolds>>-fold cross-validation on the fit slice, grouped by TimeBank-Dense document and by INLI premise so that no group spans folds. Table~\ref{tab:opt-settings} lists the optimizer settings of each configuration.
<<block:opt_settings>>

\paragraph{Split hygiene and anti-gaming.}
\texttt{optimize} loads only train and validation. The test file is hashed, never parsed, and \texttt{evaluate\_final} refuses a run whose test file changed after optimization started. Reflection sees train examples, gold labels and failures. Validation is used only for selection. A candidate program is rejected if it contains any train or validation input substring of <<block:opt.leakChars>> or more characters, any example ID, a label-leaking field, or if it grows beyond <<block:opt.sizeRule>> unless the configuration sets an explicit cap (Table~\ref{tab:opt-settings}, the numbers of the logistic regression are excluded from these checks). Demonstrations come only from train, are asserted to be train examples, and a train example is never its own demonstration. Over all test consumptions, the hygiene counters written at consumption time report <<block:opt.hygieneTestInTrainVal>> test IDs in train or validation, <<block:opt.hygieneTestChanged>> test files changed after optimization and <<block:opt.demoOutside>> demonstrations from outside train.

\paragraph{One test use per run.}
\texttt{evaluate\_final} evaluates one representative of a finished run on test, writes \texttt{test\_consumed.json} and appends a row to the study's test ledger. A second evaluation of the same run is refused. The ledger holds <<block:opt.ledgerRows>> rows (Table~\ref{tab:opt-ledger}), <<block:opt.rowsPerTask>> per task plus the LifeTox pilot. The third pre-registration contributes <<block:opt.v3LedgerSplit>> rows, and the seed reruns of its chosen configuration <<block:opt.seedLedgerSplit>> rows, one per rerun (Appendix~\ref{app:opt-seeds}). For every task all rows carry <<block:opt.testHashesPerTask>> test-file SHA-256, i.e.\ every row of a task was scored on the byte-identical file. Because the ledger allows one consumption per run, the best Jev-adapted GEPA candidate run as written was not evaluated: each Jev-adapted GEPA run's test split was consumed by <<block:opt.tailoredZeroNote>>.
<<block:opt_ledger>>

\subsection{Program IR, fitting, answers as features and demonstrations}\label{app:opt-ir}
\jevopt{} runs as an extension of the OMP agent harness plus a Python worker that talks to it over JSONL RPC (Figure~\ref{fig:arch}). OMP owns every optimizer LLM call, GEPA owns the reflective search, and \jevopt{} owns only what is Jev-specific. The worker has no LLM SDK installed, and every run's \texttt{environment.json} lists <<block:opt.forbiddenLoaded>> forbidden LLM-SDK modules loaded in the worker (\texttt{jev-opt/\allowbreak docs/\allowbreak architecture.md}):
<<block:opt_ownership>>
A \jevopt{} program is data: a YAML graph over the task's input fields whose nodes are Jev questions or routing steps and whose leaves are task labels (Table~\ref{tab:opt-ir}). Validation enforces unique node IDs, an acyclic graph in which every node and label is reachable, Jev's limits (<<block:opt.jevLimits>>) and the hygiene checks above.
<<block:opt_ir>>

\paragraph{Numeric fitting.}
Numeric parameters are never shown to the LLM. They are fit on the recorded Jev answers of the fit slice (first pre-registration: <<block:opt.v1FitExamples>> examples) by coordinate search that maximizes the task's primary metric (\texttt{jev\_opt/\allowbreak fit.py} as of <<block:opt.v3CodeTime>>): per-option additive log-probability class biases on Choice nodes (grid <<block:opt.classBiasGrid>>, decoded as $p_k e^{b_k}$ with probabilities floored at <<block:opt.featFloor>>, from the second pre-registration on), one-vs-rest thresholds (grid <<block:opt.probGrid>>) or logit biases (grid <<block:opt.biasGrid>>), branch thresholds (grid <<block:opt.probGrid>>) and Score cutoffs on the expected level (step <<block:opt.cutoffStep>> on $[0, L{-}1]$ for $L$ levels). A final temperature minimizes train log loss (bounded search on <<block:opt.tempBounds>>). It does not change the argmax. Nodes inside stage decision subtrees are frozen because changing them would change derived state.

\paragraph{Answers as features (third pre-registration).}
A Choice, one-vs-rest or Score node may carry a head that replaces its argmax by a classifier over the answers of one or more questions (``views''): Choice log-probabilities, Noul logits, and Score level probabilities and their logs, with probabilities floored at <<block:opt.featFloor>>. For classification the head is an L2 multinomial logistic regression on standardized features, $C \in$ <<block:opt.lrC>> $\times$ class weights $\in$ \{none, balanced\}, chosen by the primary metric of the out-of-fold predictions (ties: lower log loss). For QWK tasks it is a ridge regression onto the ordinal index, $\alpha \in$ <<block:opt.ridgeAlpha>>, with cutoffs fit on the out-of-fold predictions. The chosen model is refit on the whole fit slice and stored with the standardization folded into its weights. Views beyond the first are asked only as features and share the node's request when their state is identical.

\paragraph{Demonstrations (second and third pre-registration).}
A program may add labelled train examples to the Jev state of its question nodes as the reserved field \texttt{examples}: the $k$=<<block:opt.demoK>> nearest train examples by cosine similarity of <<block:opt.embedModel>> at revision <<block:opt.embedRevision>> (local, CPU, not an LLM call), taken nearest-first round-robin over labels (class-balanced). The demonstrations of an example are chosen once and kept with its result, so offline re-decoding sees the same state.

\paragraph{Structural operators and representatives.}
GEPA program search (first and second pre-registration) mixes LLM proposals with deterministic structural operators: <<block:opt_operators>>. LLM proposals may also decompose a node into stages, add derived state or regroup labels. Every proposal is validated against the IR with one retry. After a GEPA run \jevopt{} scores every fully evaluated candidate on the whole validation split and writes representatives: \texttt{top-quality} (highest validation primary metric), \texttt{best-quality} (the simplest candidate not significantly worse than GEPA's best), \texttt{best-cheap} and \texttt{best-calibrated}. GEPA with demonstrations and calibration adds per-condition \texttt{top-jev-text-opt}/\texttt{top-jev-program-opt} and ensembles \texttt{ensemble-3}/\texttt{ensemble-5}, and question optimization and Jev-adapted GEPA add the tagged representatives described below. The protocols fix which representative is tested.

\subsection{First pre-registration: the paper's question, GEPA text search and GEPA program search}\label{app:opt-v1}
The first protocol (Table~\ref{tab:opt-prereg}, recorded before the first run) fixes, per task, three independent single-condition runs, <<block:opt.nRuns>> in total: \emph{zero}, the paper's question (as a \jevopt{} program, no optimization, numeric fitting only where the IR requires it, i.e.\ the DREsS Score cutoffs), \emph{text}, GEPA text search (wording only, topology frozen), and \emph{program}, GEPA program search (wording plus Jev-specific structure). Each run consumes its test split once with its \texttt{top-quality} candidate. The primary outcome is the test difference to zero with a paired bootstrap 95\% CI. In the text condition each text slot (node instructions, option criteria, Score levels, input descriptions) is a separate GEPA component edited by GEPA's own reflective proposer. In the program condition the whole YAML is one component and \jevopt{} supplies the proposals. Reflection sees minibatches of <<block:opt.minibatch>> train examples. Candidates are scored on nested stratified validation subsets of <<block:opt.progressive>> examples and promoted to the full split only while they stay within the bootstrap interval of the current best. The optimizer model is the OMP session model (<<block:opt.optimizerSpec>>, the resolved models are listed in Appendix~\ref{app:opt-repro}).

Table~\ref{tab:opt-budget} gives the search volume. GEPA adds a proposal to the pool only if it beats its parent on the minibatch, so reflection calls far exceed candidates, particularly in the text condition, where each call edits one component. Over the <<block:opt.nRuns>> runs the optimizer made <<block:opt.gepaOptCalls>> LLM calls (<<block:opt.gepaReflCalls>> reflection, <<block:opt.gepaStructCalls>> structural) costing <<block:opt.gepaOptCost>>. Jev executed <<block:opt.gepaJevCalls>> decisions during optimization (<<block:opt.jevTokensM>>M input tokens, <<block:opt.gepaJevCost>> counting cache-served repeats). The runs' response ledgers hold <<block:opt.distinctReq>> distinct requests (<<block:opt.distinctCost>>), including <<block:opt.ledgerTestRecords>> test requests. Optimization wall-clock sums to <<block:opt.wallHours>> hours.
<<block:opt_budget>>
The execution notes appended to the protocol record two failed launches of the runner script, in which no text or program run was optimized or tested, and a fix of the run-level log-loss/Brier computation made during the study (label columns misaligned when a task's labels are not sorted). Primary metrics, GEPA's per-example objectives and the \texttt{top-quality} selection were unaffected, and the calibration numbers below are recomputed with the fixed metric.

\paragraph{Calibration.}
Brier score and ECE on test (lower is better), recomputed by \texttt{analyze.py} from the tested programs' outputs:
<<block:opt_calibration>>
Optimization targets the primary metric. Brier improves with optimization on the classification tasks, but ECE rises for TB-Dense and SMARTER. DREsS is mixed.

\paragraph{Representatives.}
Candidate indices chosen by GEPA and by each representative rule (0 = seed, parentheses give the size of the tie set of \texttt{best-quality}, and Front lists the candidates on GEPA's Pareto front over quality, calibration, calls and tokens):
<<block:opt_representatives>>
GEPA ranks by the mean per-example score it optimizes (correctness, class-balanced correctness for macro-F1, ordinal partial credit for QWK), which can differ from the primary metric. This is why \texttt{top-quality} and GEPA's choice diverge in several runs. \texttt{best-quality} often falls back to an early or seed candidate: most validation differences between candidates are within bootstrap noise. For the zero runs and TB-Dense text-opt, which started before the Brier fix, \texttt{best-calibrated} was selected with the misaligned Brier score.

\paragraph{What the GEPA-searched programs look like.}
Table~\ref{tab:opt-programs-v1} summarises structure and size.
<<block:opt_programs_v1>>For INLI, program-opt turned the flat label Choice into a two-level hierarchy: a first Choice decides whether the passage makes the hypothesis true, false or leaves it undetermined, and \emph{true} is routed to a second Choice between explicature and implicature. For SMARTER, the flat Choice became a coarse Choice over the post's rhetorical focus (in-group, out-group, explicit threat) with a fine Choice under each coarse option. For TB-Dense, the Choice became a one-vs-rest node with one Noul per relation except INCLUDES, train-fitted thresholds and INCLUDES as the fallback (introduced by the deterministic \texttt{choice\_to\_one\_vs\_rest} and \texttt{set\_abstain} operators, followed by text edits). All Nouls share one request. For EZ-STANCE, edit intent and the DREsS rubrics the tested program keeps the seed topology. The gains there come from rewritten instructions and criteria plus refitted numbers. In every DREsS program, including the zero seed, the fitted Score cutoffs contain repeated values, so some grade bands are empty and those grades are never predicted.

\subsection{Second pre-registration: GEPA with demonstrations and calibration}\label{app:opt-v2}
The second protocol asks whether class calibration, demonstrations in state, length-limited GEPA and ensembling raise Jev's test performance beyond the first pre-registration on the same tasks and splits. Its configurations, all produced on train and validation only, are: \textbf{A}, the paper's question with Choice class biases and temperature (DREsS: Score cutoffs) fit on the fit slice, no LLM; \textbf{B}, A plus <<block:opt.demoK>> class-balanced nearest train examples as demonstrations; \textbf{C}, GEPA text-opt from B with at most <<block:opt.maxComponentChars>> characters per text field and the <<block:opt.reflPool>>-example reflection pool, calibration refit per candidate; \textbf{D}, GEPA program-opt from B under the same limits, with demonstration operators enabled; \textbf{E}, the mean label distribution of the <<block:opt.ensembleK>> distinct C and D candidates by validation, with per-label biases and temperature fit on the fit slice. The better of C and D by validation was rerun with optimizer seeds 1 and 2 (representatives \texttt{top-quality} and \texttt{ensemble-3}, Appendix~\ref{app:opt-seeds}). C and D share one run per task.

\paragraph{Selection rule.}
Among all finished configurations of a task, including the seed reruns, the protocol chooses the highest validation primary metric. Ties within <<block:opt.v2Tie>> points go to fewer Jev calls per example, then to the earlier row. The generator re-applies this rule to the validation summary and checks it against the chosen configurations the protocol recorded before any test use and against the tested runs. <<block:opt.v2ProtocolDate>>. The rule is therefore fixed before any test access but not before validation results were visible.
<<block:opt_v2_chosen>>
<<block:opt_val_v2>>

\subsection{Third pre-registration: question optimization}\label{app:opt-v3}
The third protocol (recorded before its first run, one amendment before any of its validation results) asks whether a logistic regression on Jev's full answer distributions, aggressive batched rewriting of the questions, and seed programs written for Jev instead of copied from the papers raise test performance. All its runs use the fit slice and folds above and ran \jevopt{} as of <<block:opt.v3CodeTime>>.

\paragraph{Seeds written for Jev.}
Before any scoring, \texttt{/jev-seeds} asked the optimizer model, once per kind and task, for <<block:opt.seedsPerKind>> \emph{rephrased} programs (the paper's nodes, option keys and Score levels, every text rewritten following TypeSafe's authoring guidance) and <<block:opt.seedsPerKind>> \emph{decomposed} programs (any primitive or topology per routed node). The writer saw the task and label definitions, the paper seed program, the state field names and a fixed guidance text, and no data example or gold label, so every seed is itself a zero-shot program. Seeds were validated against the IR and hygiene with one retry and stored with their rationales and cost (<<block:opt.seedCalls>> calls, <<block:opt.seedCost>> in total).

\paragraph{Configurations.}
\textbf{A} is the paper's question with a logistic regression on the answers of each question node (no LLM). \textbf{P} is the optimized questions (representative \texttt{top-single}). \textbf{P+head} stacks the questions of several P variants under one logistic regression (\texttt{top-stacked}). \textbf{N}, the rephrased question, is the rephrased seed with the best zero-shot \emph{train} score, run exactly as written with no fitted number. It is a zero-shot row and was tested for every task regardless of its validation score. The rephrased seed used by N (index among the task's <<block:opt.seedsPerKind>>) was: <<block:opt_v3_native_seed>>. Question optimization works on the program's question nodes (``slots''), which in all seeds are leaves reached only through \texttt{switch} nodes:
\par{\scriptsize
\begin{verbatim}
round 0: for each slot, variants = its question in the
         paper seed (with and without <<block:opt.demoK>> kNN demos)
         + the slot's nodes of every generated seed,
         routing replaced by an LR head
ask all variants on the fit slice (one packed request
  per example and state); fit each variant's LR head
  with <<block:opt.cvFolds>>-fold CV; score = out-of-fold primary metric
for round r = 1..<<block:opt.popRounds>> (slots take turns):
    one OMP call sees the slot's <<block:opt.popElite>> best variants, a
      per-class summary and the <<block:opt.popWorst>> worst out-of-fold
      train examples; it proposes <<block:opt.popSize>> radically diverse
      variants (reframed instructions and criteria,
      Choice / one-vs-rest Nouls / Score, other state
      fields, <= <<block:opt.popAux>> aux cue questions, optional demos)
    validate (IR, hygiene, <= <<block:opt.maxComponentChars>> chars per field)
    ask all new variants in one packed request per
      example; fit + CV each variant's LR head
    stop starting rounds once fresh Jev spend >= <<block:opt.popJevCapPlain>>
programs = combinations of per-slot variants, scored
  exactly offline from out-of-fold predictions
  (exhaustive up to <<block:opt.maxCombos>>, else coordinate ascent)
single  = the <<block:opt.popValCands>> best programs by train CV,
          refit and scored on validation
stacked = per slot, forward selection of 2..<<block:opt.popStackMax>> views
          from the single programs by train CV
top-single / top-stacked = best on validation
\end{verbatim}
}
\noindent The run also scores, without further choice, every seed as written (\texttt{zero-*}), each seed kind's best seed with a head (\texttt{head-*}) and the best program descending from each kind (\texttt{search-*}), so the contribution of the seeds can be read off Table~\ref{tab:opt-val}.

\paragraph{Selection rule.}
Fixed before any validation result of this configuration: start from A. P replaces the current choice if its validation primary metric is at least <<block:opt.v3Margin>> point higher, then P+head replaces the current choice if it is at least <<block:opt.v3Margin>> point higher than that. The margin guards against the optimism of choosing P and P+head among several candidates on the same split. The generator re-applies the rule to the validation summary and checks it against the choices the protocol recorded and against the tested runs.
<<block:opt_v3_chosen>>
<<block:opt_val_v3>>
<<block:opt_v3_seeds>>

\subsection{Jev-adapted GEPA (validation study)}\label{app:opt-tailored}
A separate protocol, recorded before the first Jev-adapted GEPA run, pre-registers a validation-only study of whether GEPA adapted to Jev's decision model finds better wording than the generic GEPA text search at the same optimizer budget when both are scored after a fitted bias. It uses the fit slice, folds and rephrased seeds of question optimization, and \jevopt{} as of <<block:opt.tailoredCodeTime>>. Each task runs one GEPA text search (topology frozen) per start (<<block:opt.tailoredStarts>>, the latter being the rephrased seed with the best zero-shot train score). GEPA keeps its Pareto parent selection, progressive validation and stopping, and \jevopt{} fills five hooks: \emph{confusable-pair targeting}, \emph{minimal-pair evidence}, \emph{authoring rules}, \emph{fitted-bias scoring} and \emph{likelihood-based acceptance}:
\par{\scriptsize
\begin{verbatim}
for start in (paper, rephrased):         # budget split evenly
  GEPA text search from start, merge off; per iteration:
    parent  <- GEPA Pareto selection
    pair targeting: pair <- Choice option pair of the
        parent with the largest error mass x
        (1 - pairwise AUC) on the fit slice
        (Score-only programs: round-robin nodes)
    minimal pairs: up to <<block:opt.minimalPairs>> similar train inputs
        with gold labels on the two sides of the pair,
        with Jev's probabilities
    authoring rules: one OMP call with the fixed Jev
        rulebook rewrites both criteria + the node's
        instructions jointly
        (<= <<block:opt.maxComponentChars>> chars per field)
    fitted-bias scoring: fit class biases / cutoffs
        + temperature on the fit slice; fit-slice
        examples scored out of fold
    likelihood-based acceptance: minibatch = <<block:opt.acceptN>> train
        examples (class-balanced for macro-F1,
        label-proportional otherwise); accept child
        iff mean log P(gold) (floor <<block:opt.logpFloor>>) rises and
        the minibatch metric (after bias) drops by
        <= <<block:opt.acceptTol>> points
  stop when the start's OMP or fresh-Jev share is spent
score every candidate also as written and with a
  logistic regression
\end{verbatim}
}
\noindent The OMP budget per task is the task's GEPA program search spend and the fresh Jev budget <<block:opt.tailoredJevCap>> per run, each split evenly over the starts (Table~\ref{tab:opt-tailored-budget}). The protocol reads a Jev-adapted GEPA row as better than another only if its validation score is at least <<block:opt.tailoredMargin>> point higher. Its rows are the best fitted candidate by validation (a ``Jev + fitted numbers'' system), the best candidate run as written (zero-shot) and the best candidate with a logistic regression on its answers. All three are chosen on the validation split they are reported on (Table~\ref{tab:opt-val-tailored}). The ablations without authoring rules or likelihood-based acceptance foreseen by the protocol (``if budget allows'') were not run. The protocol does not pre-register a test round. Its best fitted candidate, fixed as the representative by the protocol before any result, was nevertheless tested once per task at the authors' request after the third pre-registration's tests (recorded <<block:opt.tailoredTestTimes>>). These rows are reported as such and are not a pre-registered test.
<<block:opt_tailored_budget>>
<<block:opt_val_tailored>>

\subsection{Test results}\label{app:opt-test}
Table~\ref{tab:opt-test-all} lists the <<block:opt.nTested>> test evaluations of the optimizer-seed-0 runs: per task, the paper's question, GEPA text search and GEPA program search, the selected configurations of GEPA with demonstrations and calibration and of question optimization, the rephrased question and Jev-adapted GEPA, each with its differences to the paper's question and to GEPA with demonstrations and calibration. The other <<block:opt.seedLedgerSplit>> test evaluations of the study, those of the chosen configuration's seed-1 and seed-2 reruns, are in Appendix~\ref{app:opt-seeds} (Table~\ref{tab:opt-seeds-headline}). Only the optimized questions are the pre-registered main result. The first pre-registration had no single pre-registered choice, and the others are references or ablations. The test analysis re-executed every representative of Table~\ref{tab:opt-test-all} offline from its run's Jev cache with network posting disabled and found <<block:opt.discrepancies>> discrepancies with the scores recorded at consumption.
<<block:opt_test_all>>
Table~\ref{tab:opt-valtest} sets the validation score of every tested representative next to its test score. For the optimized-questions rows validation exceeds test by <<block:opt.v3OptimismMean>> points on average and by up to <<block:opt.v3OptimismMax>>, while test exceeds validation on <<block:opt.v3TestAboveVal>> (over all three optimizer seeds: <<block:opt.seedsOptimism>>, Appendix~\ref{app:opt-seeds}). For the rephrased question, which involves no validation choice, the mean gap is <<block:opt.nativeGapMean>> points, so averaged over tasks the gap of the chosen rows is <<block:opt.v3VsNative>> the gap between the two splits without selection. The DREsS validation split is small, and on it the ranking of candidates did not hold on test in GEPA program search: the language rubric's program-search candidate gained on validation and lost on test (<<block:opt.dressLangV1>>). TB-Dense's validation split is the standard dev split of a few documents, and validation exceeds test there in <<block:opt.tbAllOptimistic>> rows.
<<block:opt_valtest>>

\paragraph{LifeTox pilot.}
<<block:opt.lifetoxBefore>>, a pilot run on LifeTox exercised all three conditions of the first pre-registration in one run with a smaller budget, on its own splits: train and validation from LifeTox's train split and a label-stratified test sample of <<block:opt.lifetoxTestN>> items of the released test split, not the zero-shot evaluation's test set. It was not pre-registered. Its first test evaluation was killed while queued and made no test call. At the authors' request the run's <<block:opt.lifetoxRep>> representative was later tested once.
<<block:opt_pilot>>

\subsection{Optimizer-seed variance}\label{app:opt-seeds}
\paragraph{The chosen configuration over three optimizer seeds.}
How much the optimized questions move with the optimizer seed was measured under a separate protocol (Table~\ref{tab:opt-prereg}), recorded before any seed-1 or seed-2 run existed and before any of their test evaluations. Seed 0 is the pre-registered run of Appendix~\ref{app:opt-v3}, reused and not rerun. Seeds 1 and 2 are new runs of P and P+head on all <<block:opt.nTasks>> tasks. Everything else equals seed 0: task files, seed programs, data and splits, budgets (<<block:opt.popRounds>> rounds of <<block:opt.popSize>> variants, fresh Jev spend capped at <<block:opt.popJevCap>> per run, the generator checks that each rerun's configuration differs from seed 0's only in the optimizer seed) and code (\jevopt{} <<block:opt.seedCodeNote>>). The optimizer seed changes the <<block:opt.fitExamples>>-example fit slice, the cross-validation folds of the logistic regression and the search's sampling. The OMP proposer is stochastic in any case.

The selection rule was fixed before any seed-1 or seed-2 result existed: per run, start from P. P+head replaces it iff its validation primary metric is at least <<block:opt.v3Margin>> point higher. This is the third pre-registration's rule without A (the paper's question with a head, a different method that was not rerun). Applied to seed 0 it reproduces <<block:opt.seedsReproduced>> registered choices, because A was at least <<block:opt.seedsAGap>> points below P on every task. The chosen program of each run was evaluated once on test (<<block:opt.seedLedgerSplit>> ledger rows), with no reselection after a test score was seen. The selection was recorded before the first of these test evaluations (Table~\ref{tab:opt-prereg}). Table~\ref{tab:opt-seeds-headline} lists every run and Table~\ref{tab:opt-seeds-summary} the mean and sample standard deviation over the three seeds. P+head was chosen in <<block:opt.seedsHeadChosen>> seed-1/2 runs. The test standard deviation ranges <<block:opt.seedsTestSd>> points, the validation standard deviation reaches <<block:opt.seedsValSd>>, and validation exceeds test by <<block:opt.seedsOptimism>>. <<block:opt.seedsVsAudit>>. The <<block:opt.seedsCost>>.
<<block:opt_seeds_headline>>
<<block:opt_seeds_summary>>

\paragraph{Earlier validation-only reruns: GEPA with demonstrations and calibration.}
The second protocol reran the better of C and D with optimizer seeds 1 and 2 and compared them on validation (Table~\ref{tab:opt-seeds}). The spread of the top-quality validation score across the three seeds is largest for <<block:opt.maxSpreadTask>>: <<block:opt.ezSeeds>> for seeds 0, 1 and 2, a range of <<block:opt.ezSpread>> points from identical data, settings and budget. The median range over tasks is <<block:opt.medianSpread>> points. A single GEPA run is therefore a noisy estimate of what the configuration can reach, and part of any difference between configurations on one task is optimizer seed. (The seed \emph{programs} of question optimization, each run as written, are a different notion of seed, they are in Table~\ref{tab:opt-v3-seeds}.)
<<block:opt_seeds>>

\subsection{Candidate reranking and fitted thresholds on large-label evaluated tasks}\label{app:opt-leverc}
A separate evaluated-task addendum added new test rows, never replacements, for evaluated tasks whose gap is a decoding problem over many labels. Each reuses the zero-shot evaluation's stored Jev answers wherever the request is unchanged. \textbf{TTP}: round 1 is the zero-shot evaluation's Noul scan over all ATT\&CK profiles. Round 2 asks one Choice over Jev's own top-$k$ profiles, and the shortlist is reordered by $(1-w)\,p_{\text{noul}} + w\,p_{\text{choice}}$ ($k$ and $w$ fit on dev P@1). \textbf{Wiki-WD}: the zero-shot evaluation's round-1 answers decoded with per-property log-biases (grid up to $\pm$<<block:opt.wikiBiasMax>>). Round-3 shortlist reranks were the alternatives. \textbf{MADE}: after the zero-shot evaluation's traversal, extension rounds ask every unasked code whose parent has $P > \tau_{\text{route}}$, and codes are accepted by per-code thresholds (codes with at least <<block:opt.madeMinPos>> dev positives, shared per-level thresholds otherwise) fit on validation, the paper's own threshold protocol. Dev samples come only from train, dev or validation data, with items whose text occurs in the frozen test set dropped. TTP's Expert rows use no Expert data and test transfer.

\paragraph{Registration rule.}
A configuration was registered for exactly one test run iff the paired-bootstrap 95\% CI (<<block:opt.boot>> resamples) of its held-out dev difference to the frozen zero-shot decision lay above zero. Within a task, the cheapest configuration within noise of the best was taken, and its parameters, refit on the whole dev sample, were fixed before the test run (<<block:opt.paramShas>>, SHA-256 prefixes). During the dev round the interval method changed from percentile to basic: on MADE the zero-shot decision's dev macro-F1 was <<block:opt.madeV2Pct>>, because resamples that drop a rare code's only positive report score that code as zero. The change left TTP and Wiki-WD unchanged and moved MADE's $\tau_{\text{route}}$ from <<block:opt.madeTauBest>> to <<block:opt.madeTauChosen>>, a dev difference of <<block:opt.madeTauGap>> points that is within noise under the basic interval, where the cheaper $\tau$ is chosen.
<<block:opt_leverc>>
<<block:opt_leverc_dev>>
The dev runs cost <<block:opt.leverCDevCost>>. The registered test runs added <<block:opt.leverCTestCost>> of new Jev spend. The same protocol also designed calibration and logistic regression on the zero-shot evaluation's own answers, which were registered and tested separately (Appendix~\ref{app:opt-heads}).

\subsection{Calibration and logistic regression on the zero-shot evaluation's answers}\label{app:opt-heads}
A second evaluated-task addendum, designed in the protocol of the reranking addendum and registered in its own protocol (Table~\ref{tab:opt-prereg}), fits class biases or a logistic regression on the zero-shot evaluation's own answers for <<block:opt.headsN>> evaluated test sets: ConTempo MATRES and TB-Dense, and SMARTER's HateXplain, Latent Hate and Implicit Hate sets.

\paragraph{What a head is.}
A head reads only the Choice probabilities that the frozen paper's question returned for an item. No prompt, question or model call changes, and the head replaces only the argmax. Its features are $\log\max(p, <<block:opt.headsFloor>>)$ per label. Three kinds of candidate are compared: the zero-shot argmax, per-class log-biases fit by coordinate ascent on the task metric, and an L2 multinomial logistic regression with $C \in$ <<block:opt.headsLrC>> (<<block:opt.headsBalanced>>). The candidate is chosen by <<block:opt.headsFolds>>-fold cross-validation on the fit part of a frozen dev sample, with folds grouped by document for ConTempo, and scored once on the held-out dev part.

\paragraph{Dev samples and registration rule.}
The frozen dev samples come only from each dataset's released train and dev data, with items whose text occurs in the frozen test set dropped, and the paper's question was run on them unchanged (fit / dev items: <<block:opt.headsDevSizes>>). A set was registered for exactly one test evaluation iff the chosen head's held-out dev difference to the zero-shot argmax had a paired-bootstrap 95\% CI (<<block:opt.boot>> resamples over dev items, basic interval) above zero. A registered head was refit on the whole dev sample (fit and dev parts) and frozen with its SHA-256 (prefixes: <<block:opt.headsShas>>). The registration and the frozen heads were recorded before any test output of these classifiers existed (Table~\ref{tab:opt-prereg}). All <<block:opt.headsN>> heads passed the rule, and each was evaluated once on test.

\paragraph{Results and cost.}
Table~\ref{tab:opt-heads} compares each head with the stored zero-shot test row on the same items. <<block:opt.headsTestSummary>>. A test run of such a classifier sends exactly the zero-shot run's requests, and all <<block:opt.headsTestReq>> of them were served from the stored answers of the zero-shot runs, so the test runs added no Jev spend. The dev runs cost <<block:opt.headsDevCost>>.
<<block:opt_heads>>

\paragraph{Caveats.}
<<block:opt.headsBalance>>. <<block:opt.headsWeakest>>, the weakest registration: the rule passed it, but its dev evidence is weak. The condition is ``Jev + classifier on the answers'': a trained system fit on released train and dev data, reported alongside the zero-shot evaluation's row, which it does not replace.

\subsection{Cost}\label{app:opt-cost}
Tables~\ref{tab:opt-cost}, \ref{tab:opt-cost-v2}, \ref{tab:opt-cost-v3} and~\ref{tab:opt-seeds-headline} list every run. Optimizer spend dominates: over all configurations, including seed generation and the seed reruns of the chosen configuration, the optimizer LLM cost <<block:opt.allOptCost>>, while all <<block:opt.ledgerRows>> test evaluations of the ledger (seed reruns included) together cost <<block:opt.allJevTestCost>> of Jev. ``Jev used'' counts each distinct request of a run at Jev's reported price even if the shared cache served it, so it overstates the fresh spend where runs reuse earlier answers. Fresh spend was recorded per run only from the third pre-registration on. The evaluated-task addenda (reranking and thresholds, classifiers on zero-shot answers) have their own dev runs and are listed as totals in Table~\ref{tab:opt-cost}.
<<block:opt_cost>>
<<block:opt_cost_v2>>
<<block:opt_cost_v3>>

\subsection{Program examples}\label{app:opt-programs}
Table~\ref{tab:opt-prog-summary} describes the tested program of question optimization for every task. The listings below are two of these tested programs, for EZ-STANCE (<<block:opt.ezRep>>) and SMARTER (<<block:opt.smarterRep>>), re-serialized from the stored YAML with one deterministic abridgement: each logistic regression's weight matrix \texttt{W} and intercept vector \texttt{b} are replaced by their shape. Every other field is printed as stored, including the fitted temperature, its cross-validation record and the demonstration spec. The field that holds the logistic regression is printed as \texttt{answer\_classifier}. For contrast, <<block:opt_ensemble>>.
<<block:opt_prog_summary>>
\paragraph{EZ-STANCE, optimized questions.}
<<block:opt_yaml_ez>>
\paragraph{SMARTER, optimized questions.}
<<block:opt_yaml_smarter>>

\subsection{Reproduction}\label{app:opt-repro}
\paragraph{Commands and run directories.}
Each experiment directory has a \texttt{prepare.py} (\texttt{uv run jev-opt/\allowbreak experiments/\allowbreak <exp>/\allowbreak prepare.py}) that writes the gitignored splits and a manifest with counts and SHA-256. The drivers serialize every Jev-calling command on the project-wide lock and skip finished runs. None of the dev drivers calls \texttt{/jev-evaluate-final}. Each run's single test use is \texttt{omp -p -{}-no-session "/jev-evaluate-final <run> <representative> -{}-yes"}, and \texttt{uv run scripts/\allowbreak paper\_tables.py} regenerates this appendix. Table~\ref{tab:opt-repro} maps every method to its run directories, commands and reports. It is the only place where the study's internal run and file names appear. The analysis scripts re-execute programs from the Jev cache with posting disabled. They read the run directories where they were produced (the runs may sit under different roots, passed with \texttt{-{}-v3-root}/\texttt{-{}-tailored-root}), because each root holds the Jev response cache the run used. Every response a run used is stored in its response ledger.
\begin{table*}[p]
\centering\scriptsize\setlength{\tabcolsep}{3pt}
\begin{tabular}{p{3.3cm}p{3.9cm}p{5.2cm}p{3.6cm}}
\toprule
Method & Run directories & Driver and analysis & Report \\
\midrule
Paper's question, GEPA text search, GEPA program search (first pre-registration) & \texttt{runs/\allowbreak jev-opt/\allowbreak study-<task>-\allowbreak\{zero,text,program\}} & \texttt{jev-opt/\allowbreak experiments/\allowbreak run\_all.sh} (runs + tests); \texttt{jev-opt/\allowbreak experiments/\allowbreak analyze.py} & \texttt{results/\allowbreak gepa.json}, \texttt{GEPA.md} \\
GEPA with demonstrations and calibration (second pre-registration; A, B, C+D, seed reruns) & \texttt{runs/\allowbreak jev-opt/\allowbreak v2-<task>-\allowbreak\{A,B,CD,\allowbreak <cond>-s1,\allowbreak <cond>-s2\}}, \texttt{v2-dev-summary.json} & \texttt{run\_v2\_dev.sh all}; \texttt{analyze\_v2.py} & \texttt{results/\allowbreak GEPA\_v2\_val.md} \\
Question optimization (third pre-registration; A, P) and rephrased question (N) & \texttt{runs/\allowbreak jev-opt/\allowbreak v3-<task>-\allowbreak\{A,P,N\}}, \texttt{v3-dev-summary.json}; seeds in \texttt{jev-opt/\allowbreak experiments/\allowbreak <exp>/\allowbreak seeds\_v3*/} & \texttt{run\_v3\_dev.sh seeds} (writes the seeds), then \texttt{run\_v3\_dev.sh all}; \texttt{analyze\_v3.py} (\texttt{-{}-v3-root}) & \texttt{results/\allowbreak GEPA\_v3\_val.md} \\
Optimized questions, optimizer seeds 1 and 2 (seed variance) & \texttt{runs/\allowbreak jev-opt/\allowbreak v3-<task>-\allowbreak P-s\{1,2\}} (seed 0: \texttt{v3-<task>-\allowbreak P}) & \texttt{run\_seeds.sh val}; \texttt{analyze\_seeds.py -{}-select}; \texttt{run\_seeds.sh test}; \texttt{analyze\_seeds.py} & \texttt{results/\allowbreak SEEDS.md}, \texttt{results/\allowbreak seeds.json} \\
Jev-adapted GEPA & \texttt{runs/\allowbreak jev-opt/\allowbreak tailored-<task>}, \texttt{tailored-dev-summary.json} & \texttt{run\_tailored\_dev.sh main}; \texttt{analyze\_tailored.py} (\texttt{-{}-tailored-root}) & \texttt{results/\allowbreak GEPA\_tailored\_val.md} \\
All test evaluations & \texttt{runs/\allowbreak jev-opt/\allowbreak test\_ledger.jsonl}, each run's \texttt{test\_consumed.json} & \texttt{/jev-evaluate-final}; \texttt{analyze\_test\_v3.py} & \texttt{results/\allowbreak gepa\_v3\_test.json}, \texttt{GEPA\_v3\_test.md} \\
Candidate reranking and fitted thresholds & dev: \texttt{runs/\allowbreak dev\_v3/\allowbreak \{ttp,wikiwd,made\}\_v3}; test: \texttt{runs/\allowbreak \{ttp\_v3,\allowbreak relemb\_v3/\allowbreak wikiwd,\allowbreak made\_v3\}} & \texttt{scripts/\allowbreak run\_v3\_dev.sh dev\_v3/...}; \texttt{scripts/\allowbreak fit\_v3.py \{ttp,wikiwd,made,report\}}; \texttt{uv run gauntlet run} per test set & \texttt{results/\allowbreak v3\_dev/*.json}; \texttt{data/\allowbreak frozen/\allowbreak *\_params.json} \\
Classifiers on the zero-shot evaluation's answers & dev: \texttt{runs/\allowbreak dev\_v3/\allowbreak <sample>}; test: \texttt{runs/\allowbreak \{contempo,smarter\}/\allowbreak <set>+head} & \texttt{scripts/\allowbreak run\_v3\_dev.sh dev\_v3/<sample>}; \texttt{uv run python scripts/\allowbreak fit\_v3.py heads}; \texttt{uv run gauntlet run <task>+head} per test set; \texttt{scripts/\allowbreak report\_heads.py} & \texttt{results/\allowbreak HEADS.md}, \texttt{results/\allowbreak heads.json}, \texttt{results/\allowbreak v3\_dev/\allowbreak heads.json}; \texttt{data/\allowbreak frozen/\allowbreak heads/} \\
LifeTox pilot & \texttt{runs/\allowbreak jev-opt/\allowbreak lifetox-real-1} & \texttt{jev-opt/\allowbreak experiments/\allowbreak run\_all.sh} & -- \\
\bottomrule
\end{tabular}
\caption{Reproduction map: internal run directories, drivers (under \texttt{jev-opt/\allowbreak experiments/} unless a path is given), analysis scripts (run with \texttt{uv run -{}-project jev-opt/\allowbreak python python}) and reports of every method. \texttt{<task>} is the task slug, \texttt{<exp>} its experiment directory.}
\label{tab:opt-repro}
\end{table*}

\paragraph{Pre-registration timeline.}
Table~\ref{tab:opt-prereg} gives, per configuration, when the protocol was first recorded, when the tested choice was fixed, the start of the first run and the first test consumption, from the study's git history (recorded in \texttt{scripts/\allowbreak appendix/\allowbreak opt\_history.json}) and the ledgers. The generator fails if any choice was fixed after that configuration's first test.
<<block:opt_prereg>>

\paragraph{Software and models.}
All <<block:opt.nEnvRuns>> run directories record the same environment: <<block:opt.versions>>. Jev: <<block:opt.jevModel>>. Every run's \texttt{environment.json} records only the optimizer model spec \texttt{current}. The model that answered each call is in the host's usage ledger (calls per configuration and model): <<block:opt_models>> \jevopt{} is pinned to Python 3.11: Python 3.12 changed floating-point \texttt{sum()} to compensated (Neumaier) summation, which moves Score expected levels that land exactly on a fitted cutoff. The first analysis of the first pre-registration ran under Python 3.13 and differed on the DREsS content and language text-opt runs. Under the worker's Python 3.11 all recorded test scores reproduce exactly (<<block:opt.discrepanciesV1>> discrepancies for the first pre-registration, <<block:opt.discrepancies>> over all test evaluations), and the analysis scripts assert the version.

\paragraph{Re-running the paper's question.}
The paper's-question rows are not always identical to the zero-shot evaluation's scores in Table~\ref{tab:main}. For INLI, TB-Dense and SMARTER the seed sends byte-identical Jev requests (\bodyRerunIdentical{} of \bodyRerunN{} test requests hash identically under the zero-shot evaluation's canonical request hash), but \jevopt{} keys its cache differently, so Jev was called again: \bodyRerunChangedPct\% of these items received a different probability vector (median largest per-label difference \bodyRerunMedDiff), and \bodyRerunFlips{} items (\bodyRerunFlipPct\%) changed their predicted label. For EZ-STANCE and edit intent the second routed question has a different question ID. For DREsS each request carries one rubric instead of all three and the prediction is the expected level banded by fitted cutoffs instead of the argmax level.
